(1) is Theorem \claimref{transport.existence-classification-type-c}. (2): by Lemma \claimref{transport.decomposition-commutation} (which holds for $\gamma = 0$), $B_0 = A \oplus A\nu$ with $\nu a = \hat{a}\nu$ and $\nu^2 = 0$. Then $J = A\nu$ satisfies: $J \neq 0$ ($\nu \neq 0$ and the sum is direct); $1 \notin J$ (again directness), so $J$ is proper; two-sidedness, since $A J = A A \nu \subseteq J$, $J A = A \nu A = A \hat{A} \nu \subseteq J$, and $\nu J = \nu A \nu = \hat{A} \nu^2 = 0 \subseteq J$, $J\nu \subseteq A\nu^2 = 0$; and $J^2 = A\nu A \nu = A \hat{A} \nu^2 = 0$. A simple algebra has no such ideal. The quotient: $B_0 / J = (A \oplus A\nu)/A\nu \cong A$ as algebras, the projection sending $\nu$ to $0$. (3) is Theorem \claimref{transport.rule-2-ideal-fusion-sharp-form}; a nonzero square rescales to $\pm 1$ by replacing $\nu$ with $\nu/\sqrt{|\nu^2|}$. (4) is the dimension formula $\dim \Cl = 2^{\dim \mathbb{V}'}$, degenerate or not (\claimref{std-algebra.blade-basis}). The classification statement assembles (1)--(4) with Theorem \claimref{transport.rule-2-ideal-fusion-sharp-form}(3); non-isometry of $\langle +1\rangle$ and $\langle -1 \rangle$ is Sylvester's law of inertia (\claimref{synthesis.completion-recovers-the-real-clifford-algebras}(iii)).
